% !TEX root = main.tex
% !TEX program = pdflatex
% Дополнения к основной преамбуле (preamble.tex не меняется).
% Здесь только то, что нужно для этого отчёта.

\usepackage[export]{adjustbox}   % max width / max height для скриншотов
\usepackage[font=small,labelsep=endash]{caption}
\captionsetup[figure]{justification=centering}

% \fig{имя файла в images/}{подпись}{метка}
\newcommand{\fig}[3]{%
  \begin{figure}[H]
    \centering
    \includegraphics[max width=\linewidth,max height=0.38\textheight]{images/#1}
    \caption{#2}\label{#3}
  \end{figure}}

% \code{...}: команда или имя в тексте (с переносом по строкам)
\PassOptionsToPackage{obeyspaces,spaces}{url}
\usepackage{xurl}
\urlstyle{tt}
\newcommand{\code}[1]{\path{#1}}

% Листинги терминала (fancyvrb/fvextra: корректно работают с UTF-8 и кириллицей)
\usepackage{fvextra}
\fvset{fontsize=\scriptsize,breaklines=true,breakanywhere=true,
       frame=leftline,framesep=3mm,rulecolor=\color{gray}}
% короткий блок команд внутри текста
\DefineVerbatimEnvironment{cmd}{Verbatim}{}
% \lstfile{файл в listings/}{подпись}
\newcounter{lstnum}
\newcommand{\lstfile}[2]{%
  \refstepcounter{lstnum}%
  \par\medskip\noindent\textbf{Листинг \thelstnum\ \textendash\ #2}\par\nopagebreak\smallskip
  \VerbatimInput{listings/#1}\par\medskip}
